\chapter{Statik 2}

Dieses Kapitel baut auf den folgenden Abschnitten im Skript des 2. Semesters \cite{schnidrig_physik_2026}:

\begin{minipage}[t]{.38\textwidth}
    \begin{itemize}
        \item Kraft
        \begin{itemize}
            \item Die resultierende Kraft
        \end{itemize}
        \item Moment
        \begin{itemize}
            \item Das Kräftepaar
            \item Schwerpunkt
        \end{itemize}
    \end{itemize}
\end{minipage}
\hspace{.1\textwidth}
\begin{minipage}[t]{.45\textwidth}
    \begin{itemize}
        \item Gleichgewicht eines Systems
        \begin{itemize}
            \item Zentrales Kräftesystem
            \item Parallele Kräfte
            \item Allgemeines Kräftesystem
        \end{itemize}
        \item Lagerung und Freimachen
        \begin{itemize}
            \item Das Freimachen
            \item Lagerungsarten
        \end{itemize}
    \end{itemize}
\end{minipage}

\begin{infobox}
    In der Statik - als Teilgebiet der Mechanik - wird von ruhenden und unverformbaren
    (starren) Körpern ausgegangen.
\end{infobox}

Ziel der Statik ist die Annahme der Lasten an einem
Bauteil, die Bestimmung von Lagerreaktionen und Schnittgrössen. Die Einteilung der Mechanik 
und der Statik ist in Abbildung \ref{fig:übersichtMechanikStatik} gezeigt.

\begin{figure}[H]    
    \centering          %     left botm right top
    \includegraphics[clip, trim=0cm 0cm 0cm 0cm,width=.8\textwidth]{übersichtMechanikStatik.png} %.6 stands for 0.6 x textwidth
    \caption{Einteilung Mechanik und Statik}
    \label{fig:übersichtMechanikStatik}
\end{figure}

\section{Gleichgewicht für das ebene Kräftesystem}

Ein Bauteil ist dann im Gleichgewicht, wenn sich alle an ihm angreifenden Kräfte und
Momente gegenseitig aufheben. Für das freigemachte Bauteil heisst das, dass die
resultierende Kraft und das resultierende Moment verschwinden \cite{bartsch_mechanik_2022}.

\begin{infobox}
    Ein Körper befindet sich im Gleichgewicht, wenn die Summe aller angreifenden Kräfte und
    Momente gleich Null ist.
\end{infobox}

\begin{minipage}[t]{.35\textwidth}
    \begin{equbox}
        \begin{align}
            \sum F_x &= 0 \\
            \sum F_y &= 0 \\
            \sum M &= 0
        \end{align}
    \end{equbox}
\end{minipage}
\hspace{.1\textwidth}
\begin{minipage}[t]{.55\textwidth}
    \begin{figure}[H]    
        \centering          %     left botm right top
        \includegraphics[clip, trim=0cm 0cm 0cm 0cm,width=.8\textwidth]{gleichgewicht_eines_systems.png} %.6 stands for 0.6 x textwidth
        \label{fig:gleichgewicht}
    \end{figure}
\end{minipage}

\begin{bspbox}
    Ein Generator mit Eigengewicht $G = 2000 N$ läuft mit einem Drehmoment von $M = 500 Nm$ im
    Uhrzeigersinn. Zusätzlich wirkt eine horizontale Kraft $F = 500 N$. Gesucht sind die Lagereaktionen in
    den Lagern $A$ und $B$.

    \begin{figure}[H]    
        \centering          %     left botm right top
        \includegraphics[clip, trim=0cm 0cm 0cm 0cm,width=.8\textwidth]{bsp_gleichgewicht.png} %.6 stands for 0.6 x textwidth
        \label{fig:bspgleichgewicht}
    \end{figure}

    Beim Aufstellen der 3 Gleichgewichtsbedingungen sind immer verschiedene Wege möglich. Mit einer
    Momentenbedingung bezüglich eines günstigen Pols kann man bereits eine Unbekannte bestimmen.
    Der Pol muss dabei nicht innerhalb des Bauteils sein. Hier wird der Punkt $C$ gewählt, durch welchen
    die beiden Kräfte $A_y$ und $B$ gehen.

    Die Wirkungslinie der Kraft $B$ ist durch die Verbindung der Punkte $B$ und $C$ gegeben (Pendelstütze).
    Der Winkel ergibt sich aus der Geometrie wie folgt:

    \begin{equation*}
        \tan^{-1}(\frac{70mm}{140mm}) = 26.57^{\circ}
    \end{equation*}

    \begin{align*}
        \sum M_C = 
    \end{align*}

\end{bspbox}

Das gezeigte Beispiel ist ein allgemeines Kräftesystem. 
Für die Spezialfälle fallen gegebenenfalls Gleichungen weg:

\begin{itemize}
    \item Zentrales Kräftesystem
    \begin{itemize}
        \item $\sum M = 0$ kann nicht aufgestellt werden
    \end{itemize}
    \item Parallele Kräftepaar
    \begin{itemize}
        \item $\sum F$ in eine Richtung wird nicht aufgestellt
    \end{itemize}
\end{itemize}

Ausserdem kann anstelle von $\sum F_x$ oder $\sum F_y$ eine weitere Momentengleichung $\sum M_{II}$ an 
einem anderen Punkt eingeführt werden.